\section{Introduction}
\label{sec:intro}

In modern deep learning, the paradigm of pre-training a foundation model and fine-tuning it on specific downstream tasks has achieved remarkable success. However, deploying a separate fine-tuned model for each target task in production is highly impractical. The storage, memory, and computational costs scale linearly with the number of tasks, creating a prohibitive bottleneck for resource-constrained edge devices, real-time embedded systems, and large-scale cloud applications. 

To address this challenge, \textbf{parameter-space model merging} has emerged as an attractive, zero-overhead alternative \cite{wortsman2022model, matena2022merging}. Techniques such as Task Arithmetic \cite{ilharco2022editing}, TIES-Merging \cite{yadav2023ties}, and AdaMerging \cite{adamerging} fuse multiple fine-tuned models directly in their parameter space, producing a single merged model that can perform all target tasks simultaneously. Because this fusion is performed offline, the resulting merged model incurs absolutely zero additional inference latency or memory footprint during deployment.

\begin{figure}[t]
\centering
\includegraphics[width=\columnwidth]{latency_vs_accuracy.png}
\caption{\textbf{Pareto Frontier of Hybrid-Router.} Our proposed Hybrid-Router (Crimson) sweeps different dynamic partition depths $k \in \{0, 1, 2, 4, 12, 14\}$ on a Vision Transformer (ViT-Tiny, $L=14$ layer groups). Our framework forms a highly favorable Pareto frontier: shifting from fully dynamic ($k=14$) to hybrid ($k=4$) reduces wall-clock parameter blending time from 10.28 ms to 2.95 ms (a \textbf{71.3\%} speedup) and saves 71.4\% of task-vector storage, while maintaining a high joint accuracy of \textbf{76.75\%} (\textbf{+4.44\%} absolute gain over SOTA static AdaMerging). Peak accuracy is achieved at $k=12$ with a \textbf{+0.22\%} gain over fully dynamic ensembling while saving \textbf{14.3\%} in assembly latency.}
\label{fig:pareto}
\end{figure}

Despite these benefits, static model merging techniques suffer from a major limitation: \textbf{representation conflicts and destructive interference}. When target tasks are highly divergent or high-conflict (e.g., combining SVHN digits with CIFAR-10 natural objects), their respective task vectors point in opposing directions in the high-dimensional parameter space. Static blending averages out these directional signals, leading to catastrophic domain sacrifice where the performance on some or all tasks degrades significantly.

To bypass representation conflicts, researchers have recently explored \textbf{dynamic model merging} or test-time parameter-space routing \cite{qwsmerge, polymerge}. These methods extract features from the current input batch at runtime and use a lightweight routing head to dynamically compute scaling coefficients. These coefficients are then used to reconstruct a customized set of weights in parameter space on-the-fly before each forward pass. 

\paragraph{The Pragmatist's Real-World Bottleneck.}
While dynamic parameter-space ensembling achieves impressive multi-task accuracy, its real-world deployability is severely limited. Performing high-dimensional linear combinations of entire weight matrices (which contain tens or hundreds of millions of parameters) at runtime introduces a massive \textbf{memory-bandwidth and computational latency bottleneck}. 
For instance, in a Vision Transformer (such as a standard 14-layer ViT-Tiny backbone), dynamically ensembling all layers on-the-fly takes $10.28$ milliseconds of wall-clock time per forward pass. This memory transfer overhead completely destroys the throughput advantages of GPU/CPU hardware and is unacceptable in latency-sensitive, real-world deployment pipelines.

\paragraph{Our Proposed Solution: Hybrid-Router.}
In this work, we present a highly practical, low-latency framework called \textbf{Hybrid-Router} that bridges the gap between zero-overhead static merging and highly expressive dynamic routing. Our framework is grounded in a key architectural observation: \emph{early layers in deep networks act as task-agnostic feature extractors (learning edges, textures, and basic shapes), whereas late layers capture task-specific specialized representations.} 

This layer-wise division of labor allows us to partition the network:
\begin{enumerate}
    \item We \textbf{statically merge} the task-agnostic early layers ($1 \dots L-k$) offline using standard Uniform Merging, requiring absolutely zero computational or memory overhead during deployment.
    \item We \textbf{dynamically merge} only the final $k$ layers online using standard Softmax routing (for peak absolute accuracy) or an alternative, Softmax-free sigmoidal projection engine (\textbf{BSigmoid-Router}) which we explore to analyze the trade-offs and scaling dynamics of uncoupled task coefficients.
\end{enumerate}

\paragraph{A Highly Favorable Trade-off of Speed, VRAM, and Accuracy.}
Through systematic evaluations conducted within a controlled Parameter-Space Representation Sandbox modeling a Vision Transformer (ViT-Tiny, $L=14$) over high-conflict vision domains (MNIST, FashionMNIST, CIFAR-10, SVHN), we find that layer-wise partitioning serves as a highly effective mechanism for trading off minor, bounded accuracy for substantial, hardware-critical systems savings. 
At $k=4$, our Hybrid-Router achieves an outstanding joint mean of \textbf{76.75\%} (a massive \textbf{+4.44\%} absolute accuracy improvement over the strong offline static AdaMerging baseline of 72.31\%) while cutting weight assembly latency and task-vector Video Random Access Memory (VRAM) footprint by \textbf{71.3\%} (2.95 ms vs 10.28 ms) and discarding 71.4\% of active dynamic task vectors from memory.

Furthermore, we observe a highly intriguing phenomenon under extremely low-resource calibration constraints (e.g., a tiny 64-sample calibration budget): fully dynamic ensembling ($k=14$) can suffer from a mild form of representation overfitting due to excessive degrees of freedom in the optimization space. In contrast, our layer-wise partitioning can act as a form of \textbf{structural regularization}. Freezing early layers to a static uniform average restricts the learnable search space of the routing optimizer, preventing over-parameterized overfitting. At $k=12$, this structural constraint yields a joint accuracy of \textbf{84.79\%} (a minor \textbf{+0.22\%} gain over the fully dynamic $k=14$ baseline of \textbf{84.57\%}) while still saving \textbf{14.3\%} in assembly latency (8.81 ms vs 10.28 ms). 

While this mild overfitting paradox is primarily observed under specific sandbox constraints and represents a minor absolute accuracy delta, the broader systems-level Pareto frontier (Figure~\ref{fig:pareto}) remains spectacular and robust: reducing the active dynamic parameter space dramatically accelerates ensembling throughput while preserving or enhancing multi-task capabilities.

\paragraph{Parameter-Space Representation Sandbox.}
To evaluate these complex layer-wise and parameter-space routing dynamics with maximum reproducibility, we formulate a self-contained, mathematically rigorous, and highly accessible PyTorch-based \emph{Parameter-Space Representation Sandbox}. This proxy environment maps high-conflict image domains to a semi-orthogonal feature space, simulating the hierarchical layers of a Vision Transformer and physical weight interpolation. This provides a clean, deterministic benchmark for isolating weight ensembling algorithms.

\paragraph{Core Contributions.} Our primary contributions are:
\begin{itemize}
    \item We analyze the runtime latency and memory-bandwidth bottlenecks of test-time dynamic model merging, bringing a crucial pragmatist focus and hardware-aware metrics (VRAM, throughput) to parameter-space routing.
    \item We propose \textbf{Hybrid-Router}, which partitions deep models layer-wise to statically merge general-purpose early layers offline (enabling 71.4\% task-vector VRAM savings at $k=4$) and dynamically route only the final $k$ task-specific layers.
    \item We demonstrate that layer-wise partitioning acts as a highly favorable trade-off mechanism, yielding a superior Pareto frontier that significantly outperforms SOTA static model merging.
    \item We present an exploratory study on a Softmax-free, independent sigmoidal routing projection (\textbf{BSigmoid-Router}) to understand the mathematical and empirical trade-offs of uncoupled task coefficients, candidly analyzing its performance limits, identifying its substantial performance gap compared to standard Softmax, and demonstrating how this gap is driven by conservative scaling ceilings.
    \item We validate our model under real-world streaming noise and small batch sizes ($B=1$), demonstrating high sample-efficiency and temporal stability in the wild.
\end{itemize}
